\chapter{The Local-First Approach}
\label{ch:localfirst}

\textcite{kleppmann2019localfirst} listed what software must do if users are to
own their data: be fast, work across devices and offline, support
collaboration, last, be private by default, and leave the user in control. The
paper predates the current AI boom, but its questions also help us assess
how we use models and who controls the data.

\section{The case for local}

\textbf{Nothing leaves the machine.} Unpublished records, locations of
threatened species, embargoed sequences, draft manuscripts and interview data
stay where they are. This is not idle caution: shared chatbot conversations have
turned up in public search results, API keys included
\parencite{ibtimes2026claudechats}, and data brokers routinely ignore privacy
law \parencite{hai2026privacylaws}. For European institutions, GDPR and
data-sharing agreements are far simpler when the answer to ``where did the data
go?'' is ``nowhere''.

\textbf{No cost per query.} Once the hardware is paid for, you stop rationing:
you run the classifier over the whole corpus, not a sample, and leave jobs
running overnight.

\textbf{It works offline} --- on research vessels, at field stations, and on
institutional networks that block outside traffic.

\textbf{It lasts.} A provider retires an API on its own schedule; a model file
on your disk does not change. Research infrastructure has a ten-year horizon.

\section{Digital sovereignty and data protection}

\textcite{jrc2025opensovereignty} define EU digital sovereignty as the ability to
act independently in the digital domain while staying open --- not isolation,
but not being coerced. \textcite{anand2026digitalsovereignty} extract the
concrete goals from EU law. \textcite{hai2026sovereignty} are sceptical of the
``sovereign AI'' now on sale: much of it is the same proprietary system hosted in
a different region, which moves the data without changing who controls the
model. \textcite{esposito2026sovereign} argue that sovereignty must be designed
into a system rather than ticked off at procurement.

Governments now treat biological data the same way. Three US bills, endorsed by
OpenAI and summarised by \textcite{turner2026databills}, propose changes to how
biological data is managed. The national standards agency would define
``AI-ready'' biological datasets and set minimum requirements for federally
funded research. A single gateway would provide access to federal biological
data, with tiers based on sensitivity. The bills would also fund measurement
standards, including for biosafety and biosecurity. The head of
the commission behind them says ``AI-ready biological data must be treated like a
strategic national resource''. Turner notes that the bill does not say how a
dataset is assigned to a tier. Two lessons carry over. Data one side calls
AI-ready, another calls strategic, so control of access is a sovereignty
question. And tiered access suits sensitive ecological data, such as the
locations of threatened species --- once the rules for the tiers are written
down.

For a public institution, these arguments lead to practical questions. Could you carry on
if a provider changed its terms, raised prices or left your market? Can you
audit what you depend on? Will your infrastructure outlive the contract?

\section{The costs}

\textbf{Capability.} Frontier hosted models are better at hard reasoning, long
documents and changes spanning many files. The gap is narrowing:
\textcite{willison2026qwen38} ran Qwen 3.8 27B, an Apache-2.0 model that fits in
a 17\,GB file, offline on a laptop and on a DGX Spark. It wrote working code,
drove a coding agent through a repository, and drew accurate boxes around birds
in a photograph --- in his judgement, what the best proprietary models could do a
year earlier. The Singapore Consensus puts open-weight models three to twelve
months behind the frontier on common benchmarks \parencite{casper2026singapore}.
On environmental geospatial agent tasks, an open-weight model reached 93 per
cent of the best model's score at less than a tenth of its cost
\parencite{diazireland2026geonature}.
Not everyone thinks the gap is worth accepting. In a 2025 talk,
\textcite{wickham2026nobs} judged that for someone starting out, local models are
by and large not worth it. His advice was to begin with a hosted model, which is
more likely to give good results. A taxonomic test supports the caution. When
morphological characters were extracted from Australian daisy descriptions, a
locally run open model was considerably less reproducible than a hosted one, with
unit errors and skipped characters \parencite{schmidtlebuhn2025morphology}.

\textbf{Speed.} The same model produced 15--30 tokens per second, slow enough
that he still preferred hosted models for interactive work. For overnight batch
jobs it does not matter.

\textbf{Hardware and upkeep.} Money up front (\Cref{ch:hardware}) and time
spent on drivers, updates and debugging. In a small group that time comes out of
someone's research, and it is usually left out of the comparison.

\section{Open, closed, and in between}

\begin{table}[htbp]
\centering
\caption{Ways to use a model, and their tradeoffs.}
\label{tab:deployment}
\small
\begin{tabularx}{\textwidth}{lXX}
\toprule
\textbf{Model} & \textbf{You get} & \textbf{You give up} \\
\midrule
Hosted chat (SaaS) & No setup, best capability & Data control, versions, offline use \\
API & Programmability, scale & The same, plus cost per token \\
Managed open weights (MaaS) & Open models, no upkeep & Data still leaves; vendor dependence \\
Self-hosted open weights & Control, privacy, offline & Capability gap, hardware, upkeep \\
\bottomrule
\end{tabularx}
\end{table}

\term{Model as a Service} is the middle path: firms including Red Hat and Ollama
offer supported deployment of open-weight models, and \textcite{redhat2026hybrid}
argue for placing each workload where its needs put it. In my view, this is a
useful approach for most institutions. A group policy can make these choices
explicit:

\begin{itemize}
  \item \textbf{Always local:} unpublished data, sensitive localities
    \parencite{tulloch2018sensitive,fox2025safeguarding}, personal
    data, anything under embargo or a restrictive agreement.
  \item \textbf{Local by default:} bulk processing, extraction and
    classification over large collections.
  \item \textbf{Hosted is fine:} public data, published literature, general
    explanation, and hard reasoning where the capability gap matters.
  \item \textbf{Hosted with care:} code generation --- check what context is
    sent and whether your licences and policy allow it.
\end{itemize}

\section{A reference architecture}

Red Hat's strategy ebook by \textcite{ferris2025adaptable} makes a useful
organisational argument: agree on the outcome, invest in staff skills, understand
where data and applications reside, and keep infrastructure adaptable. It is
vendor guidance rather than an independent evaluation. For a biodiversity group,
I would translate it into a specific acceptance test before choosing a platform:
how many verified records can the workflow produce within the available staff
time and computing budget? Portability matters when a model changes, but so
does having someone able to maintain and evaluate its replacement.

The architecture follows the same principle: keep the parts that must last
simple, and allow the others to be replaced. Each layer has a specific role.

\begin{description}
  \item[Interface] Open WebUI or LibreChat in a browser, with the model
    switchable per conversation. It stores no long-term knowledge.
  \item[Routing] An OpenAI-compatible API layer, behind which a local model or a
    hosted one can be chosen per request.
  \item[Runtime] Ollama for one user, vLLM for a group. Input in, output out.
  \item[Retrieval] PostgreSQL with the \code{pgvector} extension, plus PostGIS
    for spatial data. Embeddings from \code{bge-m3}, reranking with
    \code{bge-reranker}, and LlamaIndex, Haystack or, for branching workflows,
    LangGraph to connect them.
  \item[Archive] Plain Markdown files, append-only. This is the real long-term
    memory, synced and backed up under 3-2-1 (\Cref{ch:organisation}).
\end{description}

\emph{The interface is for talking, the models are for thinking, and plain text
is for remembering.} A small proxy between interface and runtime that writes
every exchange to a dated Markdown file gives you the provenance record
\Cref{ch:reproducibility} asks for. None of this needs an account, a contract or
a network connection; \textcite{confidence2026localstack} walk through the
practical setup.

\section{A worked example: KYBELE}

KYBELE, a project the author works on, is a literature assistant for Collembola
(springtails) built on the SIB Literature Services (SIBiLS). SIBiLS also reuses
the taxonomic treatments that Plazi extracts from the literature
\parencite{agosti2022treatments}. KYBELE adds two layers
to the architecture above: an \term{orchestrator} between interface and model
that holds all the domain rules, and a reverse proxy (Caddy) handling encryption
and rate limits. A public website, Open WebUI for users signing in through
ELIXIR, and an agent skill for tools such as Claude Code all call the same
orchestrator, so the rules are written once. The model --- Qwen3-4B by default,
served by vLLM on a 16\,GB GPU and never exposed to the internet --- is the least
specialised part. The expertise sits elsewhere:

\begin{enumerate}
  \item \textbf{Vocabulary.} Questions are expanded with curated Collembola
    terms --- synonyms, accepted names, character terms --- by plain lookup.
    General models are weakest here: across ten taxonomies they handled
    specialised domains and the most specific ranks far worse than common
    ones \parencite{sun2024taxoglimpse}.
  \item \textbf{Scope in code.} Results are filtered to Collembola-annotated
    items. The boundary is a filter, not a prompt instruction the model can
    ignore.
  \item \textbf{Saying no.} If the curated evidence does not cover the question,
    the system says so instead of answering.
  \item \textbf{Citation check.} A citation that does not match a retrieved item
    is dropped or flagged.
\end{enumerate}

React-to-Me \parencite{mohammadi2025reacttome} has tested the same design on the
expert-curated Reactome pathway database. It combines keyword and vector search
and confines answers to Reactome, with links to each entry. When Reactome has
no answer, it uses trusted outside sources instead of guesses. Combined search beat
vector search alone, and blinded experts more often rated the grounded answers
higher, with significant gains in accuracy and biological detail.
\textcite{iyer2025conservationevidence} found the same over the Conservation
Evidence database. On exam-style questions about conservation interventions,
combined keyword and vector retrieval performed best by a substantial margin,
while general models used out of the box were likely to misinform. At the time of
writing KYBELE's serving layer runs and the orchestrator is designed but not
built; the design is the point here.

\subsection{Testing it}

A single accuracy score hides most of what matters. Three examples show other
ways to test.

\textbf{Test several things at once.} Karenina, from EMBL-EBI, lets experts
say in plain language what a sound answer needs and what behaviour to test, and turns that into runnable checks \parencite{carli2026karenina}. On
144 drug-discovery questions, giving seven models live access to the Open
Targets database raised accuracy from 51 to 81 per cent. However, 11 per cent
of the answers given with that access were correct but unsupported by the
retrieved evidence. Invented PubMed identifiers gave way to real papers cited
for claims they do not make. Tool use also cost 24 to 121 times more tokens.
When a user confidently
pushed back with false information, models abandoned correct answers 57 per cent
of the time --- sycophancy, measured (\Cref{ch:powers}). A second model acting
as a guard flagged 97 per cent of those reversals. On an agent benchmark,
switching the harness alone moved one model from 50 to 65 per cent, and reading
the agents' saved code separated flawed biology from failed computing.

\textbf{Ask what a benchmark can prove.} \textcite{saezrodriguez2026benchmarking}
argue that foundation models are hard to benchmark because they are meant to
embody the processes that produced the data. They ask whether such models can be
refuted, verified, or only judged by usefulness, and what part the scientific
community should play in testing them.

\textbf{Test against the previous version.} When EMBL-EBI rewrote InterProScan,
its protein annotation pipeline, as a Nextflow workflow, the authors compared
the new version's annotations of all Swiss-Prot proteins with the old one's. They
found near-identical precision and sensitivity, at about twice the speed on large
proteomes \parencite{blum2026interproscan}.

For KYBELE this becomes a test suite, rerun after every change of prompt, model
or retrieval:

\begin{itemize}
  \item a gold set of expert Collembola questions with expected sources --- did
    it find them?
  \item a grounding check --- is each claim supported by what was retrieved?
  \item a citation audit --- does the cited paper say this?
  \item a pushback test --- does it hold a correct answer when challenged?
  \item the token cost of each answer;
  \item a comparison with the previous version, so that replacing Qwen3-4B with
    another model is accepted only if the gold-set results hold.
\end{itemize}

The ecological examples in \Cref{ch:applications} suggest a practical extension
to this test suite. The following is my proposed procedure for comparing
extraction, classification and study-retrieval workflows, drawing on their
different successes and failures
\parencite{gougherty2024extraction,domazetoski2025traits,dorm2026ecologicalknowledge,munch2026phenotypes,venetsianou2026relatedness}.

\begin{enumerate}
  \item \textbf{Fix the task and the reference answers.} Specify the fields,
    trait definitions, units and acceptable evidence. Have experts independently
    annotate a sample, resolve disagreements and retain uncertain cases as
    uncertain. Set aside a final test set before adjusting prompts or thresholds.
    Keep descriptions of the same species, or samples from the same study,
    together when testing transfer to new species or studies.
  \item \textbf{Compare like with like.} Run keyword matching or a simple
    classifier alongside the language model. Include the established domain
    method where one exists. Test unaided answers, answers given a fixed source
    passage, and answers using retrieval separately. Record which documents
    each system could access; these conditions test different abilities.
  \item \textbf{Score the errors that matter.} For extraction, report precision
    and recall for each field, plus how often missing information becomes an
    invented value. Check numerical errors after converting units, stating any
    transformation used. For traits, show each class's results and balanced
    accuracy, so common classes cannot conceal failure on rare ones.
    Count malformed outputs, failures and unanswered cases across the full test
    set, even when calculating accuracy only on valid answers.
  \item \textbf{Test uncertainty and unfamiliar cases.} Include absent fields,
    rare taxa, invented names and difficult negative examples. Report
    \term{coverage}, the proportion of cases receiving an answer, alongside
    accuracy. If confidence controls which answers are accepted, check whether
    higher confidence actually predicts fewer errors. Measure the errors and
    coverage again at each proposed threshold.
  \item \textbf{Evaluate retrieval and evidence separately.} For literature
    assistants, measure how many expert-selected relevant papers appear in the
    first few results. Then check whether each generated claim is supported by
    its cited passage. For microbiome comparisons, add independent expert
    judgements of study compatibility. Agreement between text and biological
    profiles is useful corroboration, but does not substitute for those checks.
  \item \textbf{Measure variation and the whole cost.} Repeat runs and modest
    prompt changes, checking agreement on individual records. Report results
    by taxon, biome or source collection, with sample sizes and uncertainty
    intervals. Record runtime, token use, hardware and human correction time.
    Judge a replacement by verified output per unit of work and by unacceptable
    errors, with the acceptance criteria fixed before comparing models.
\end{enumerate}

In the microbiome relatedness study, the authors also tested the similarity
measures themselves by reducing sequencing depth and asking whether profiles
still retrieved their original counterparts \parencite{venetsianou2026relatedness}.
This is a useful distinction: a benchmark's comparison method needs checking
as well as the language model. Repeating a judgement consistently establishes
repeatability; an independent reference is needed to assess correctness.

\subsection{Where fine-tuning fits}

\textcite{asai2026openscholar} show the next step. OpenScholar retrieves from 45
million open-access papers and answers with citations. Without retrieval, GPT-4o
fabricated citations to recent literature in 78--90 per cent of cases;
OpenScholar's citation accuracy matched human experts'. To train a small model,
they ran their pipeline with a larger open model (Llama 3.1 70B).
They kept answers that passed a quality filter and mixed them half and half
with general instruction data. They then fine-tuned Llama 3.1 8B on the
resulting 130{,}000 examples. The 8B model beat GPT-4o by 6.1 per cent on correctness in multi-paper
synthesis, and everything is public. Retrieval still sometimes misses the most
relevant papers, and the 8B model still makes factual errors.

KYBELE plans to adapt its model with \term{LoRA}, low-rank adaptation, which
trains a small set of added weights rather than the whole model
\parencite{hu2021lora}. The same recipe fits:
the pipeline produces grounded answers to train on, and the test suite shows
whether the tuned model is actually better.

\begin{keypoint}
Put domain expertise into data, code and tests --- vocabulary, scope filter,
gold set --- before putting it into model weights. Fine-tune last, and only keep
a tuned model that passes the tests.
\end{keypoint}

\section{Decentralised ecology}

\textcite{placitelli2026seaturtle} built a sea-turtle detection system that runs
entirely in the browser of the device that took the photograph: no upload, no
server, no connection needed. This offers several advantages for citizen science.
Volunteers keep their data, the system works in the field, and it can continue
without a server budget when the grant ends. \textcite{haas2025corda} offer a programming model for
such applications, in which each device works locally and the system reconciles
the results, and \textcite{kleppmann2024bluesky} show federated identity working
at social-network scale. Biodiversity monitoring --- distributed, often offline,
spread across institutions --- is close to the ideal use.
